\section{MOPD Learning Objective}
\label{sec:setting}

We follow the notation and loss definitions of
\citet{ma2026mopdmultiteacheronpolicydistillation}, Section~3.2.
The student policy is $\pi_\theta$ and the frozen teacher for domain $d$
is $\pi_{\phi_d}$. A prompt $x$ is drawn from the training mixture and
the student generates a response $y\sim\pi_\theta(\cdot\mid x)$.
At position $t$, both models condition on the same prefix $(x,y_{<t})$.
Following MOPD, we abbreviate these conditional probabilities as
$\pi_\theta(v)$ and $\pi_{\phi_d}(v)$, where $v$ is a vocabulary token.
Teachers and student share the vocabulary and tokenizer.

\paragraph{Reverse-KL objective.}
MOPD minimizes the student--teacher reverse KL at each response position
and averages over the response length $|y|$:
\begin{equation}
\mathcal L_{\mathrm{rev\text{-}KL}}
=\mathbb E_{x,\,y\sim\pi_\theta}
\left[\frac{1}{|y|}\sum_{t=1}^{|y|}\sum_v
\pi_\theta(v)\log\frac{\pi_\theta(v)}{\pi_{\phi_d}(v)}\right].
\label{eq:reference_reverse_kl}
\end{equation}
Each prompt is scored by its assigned teacher. The factor $1/|y|$ gives
each response a mean token loss; it does not give every token in the
entire training batch equal weight. This distinction motivates
Section~\ref{sec:normalization}.

\paragraph{Sampled-token policy gradient.}
MOPD assigns each generated token a detached log-probability advantage,
$\hat A_{\mathrm{MOPD},t}=\sg[\log\pi_{\phi_d}(y_t)-\log\pi_\theta(y_t)]$,
where $\sg$ stops gradients. The clipped value is
$\hat A^{\mathrm{clip}}_{\mathrm{MOPD},t}
=\operatorname{clip}(\hat A_{\mathrm{MOPD},t},-A_{\max},A_{\max})$,
and the policy-gradient (PG) loss is
\begin{equation}
\mathcal L^{\mathrm{PG}}_{\mathrm{MOPD}}(\theta)
=-\mathbb E_{x,\,y\sim\pi_\theta}
\left[\frac{1}{|y|}\sum_{t=1}^{|y|}
\hat A^{\mathrm{clip}}_{\mathrm{MOPD},t}\log\pi_\theta(y_t)\right].
\label{eq:mopd_pg}
\end{equation}
This loss uses the sampled token $y_t$ at each position. Clipping is part
of the practical recipe and must be recorded when comparing losses.

\paragraph{Teacher top-$k$ distillation.}
Let $\mathcal T_t^d=\operatorname{TopK}_k(\pi_{\phi_d}(\cdot\mid x,y_{<t}))$
be the teacher's $k$ most probable tokens. MOPD instead evaluates
\begin{equation}
\mathcal L^{\mathrm{TopK}}_{\mathrm{MOPD}}(\theta)
=\mathbb E_{x,\,y\sim\pi_\theta}
\left[\frac{1}{|y|}\sum_{t=1}^{|y|}\sum_{v\in\mathcal T_t^d}
\left(\pi_\theta(v)\log\frac{\pi_\theta(v)}{\pi_{\phi_d}(v)}
-\pi_\theta(v)+\pi_{\phi_d}(v)\right)\right].
\label{eq:mopd_topk}
\end{equation}
These are full-vocabulary probabilities, without renormalization over
the selected tokens. The correction $-\pi_\theta(v)+\pi_{\phi_d}(v)$
makes each selected-token term minimal when the two probabilities match;
it does not recover the omitted vocabulary terms. With the entire
vocabulary included, the correction sums to zero and the loss reduces
to Equation~\ref{eq:reference_reverse_kl}.

\paragraph{Gradients measured in this paper.}
We differentiate the chosen loss on fixed student-generated responses,
holding prefixes, masks, teachers, and detached advantages fixed.
These local training gradients do not include derivatives through the
generation of the responses or their lengths. The unclipped sampled-token
gradient matches the full reverse-KL gradient in expectation at a fixed
prefix; clipping, top-$k$ truncation, and trajectory-dependent weighting
need separate treatment (Section~\ref{sec:density}).
